\begin{figure}
    \centering
    \begin{subfigure}{0.5\textwidth}
        \includegraphics[width=\linewidth]{nt_RIO_0.027.png}
        \caption{}
        \label{fig:rio_0.027_nt}
    \end{subfigure}%
    \begin{subfigure}{0.5\textwidth}
        \includegraphics[width=\linewidth]{scaled_nt_RIO_0.027.png}
        \caption{}
        \label{fig:rio_0.027_scaled_nt}
    \end{subfigure}
    
    \begin{subfigure}{0.5\textwidth}
        \includegraphics[width=\linewidth]{nr_RIO_0.027.png}
        \caption{}
        \label{fig:rio_0.027_nr}
    \end{subfigure}%
    \begin{subfigure}{0.5\textwidth}
        \includegraphics[width=\linewidth]{gamma_tau_RIO_0.027.png}
        \caption{}
        \label{fig:rio_0.027_gamma_tau}
    \end{subfigure}
    \caption{Here are represented the important features about the phase transition in Rio De Janeiro for UCI = 0.027.
    In \ref{fig:rio_0.027_nt} is represented the number of people trapped in the road network at time t. In \ref{fig:rio_0.027_scaled_nt}, the vertical line representing $t_{\infty}$. we see the values $n(t)t^{\alpha}$ over $\epsilon (t-t_0)$, with $\epsilon = \frac{R^2 - R_c^2}{R_c^2}$ with t up to $t_{start} + t_{end_{fit}}$ In \ref{fig:rio_0.027_nr} showing the phase transition requirement. We plot n(R,$t_{\infty}$) with $\beta = $ 0.144. In \ref{fig:rio_0.027_gamma_tau} we see the linear increase of $\tau$ compared to $\Gamma$}
    \label{fig:phase_transition_rio_0.027}
\end{figure}

\begin{figure}
    \centering
    \begin{subfigure}{0.5\textwidth}
        \includegraphics[width=\linewidth]{nt_RIO_0.091.png}
        \caption{}
        \label{fig:rio_0.091_nt}
    \end{subfigure}%
    \begin{subfigure}{0.5\textwidth}
        \includegraphics[width=\linewidth]{scaled_nt_RIO_0.091.png}
        \caption{}
        \label{fig:rio_0.091_scaled_nt}
    \end{subfigure}
    
    \begin{subfigure}{0.5\textwidth}
        \includegraphics[width=\linewidth]{nr_RIO_0.091.png}
        \caption{}
        \label{fig:rio_0.091_nr}
    \end{subfigure}%
    \begin{subfigure}{0.5\textwidth}
        \includegraphics[width=\linewidth]{gamma_tau_RIO_0.091.png}
        \caption{}
        \label{fig:rio_0.091_gamma_tau}
    \end{subfigure}
    \caption{Here are represented the important features about the phase transition in Rio De Janeiro for UCI = 0.091.
    In \ref{fig:rio_0.091_nt} is represented the number of people trapped in the road network at time t. In \ref{fig:rio_0.091_scaled_nt}, the vertical line representing $t_{\infty}$. we see the values $n(t)t^{\alpha}$ over $\epsilon (t-t_0)$, with $\epsilon = \frac{R^2 - R_c^2}{R_c^2}$ with t up to $t_{start} + t_{end_{fit}}$ In \ref{fig:rio_0.091_nr} showing the phase transition requirement. We plot n(R,$t_{\infty}$) with $\beta = $ 0.092. In \ref{fig:rio_0.091_gamma_tau} we see the linear increase of $\tau$ compared to $\Gamma$}
    \label{fig:phase_transition_rio_0.091}
\end{figure}

\begin{figure}
    \centering
    \begin{subfigure}{0.5\textwidth}
        \includegraphics[width=\linewidth]{nt_RIO_0.092.png}
        \caption{}
        \label{fig:rio_0.092_nt}
    \end{subfigure}%
    \begin{subfigure}{0.5\textwidth}
        \includegraphics[width=\linewidth]{scaled_nt_RIO_0.092.png}
        \caption{}
        \label{fig:rio_0.092_scaled_nt}
    \end{subfigure}
    
    \begin{subfigure}{0.5\textwidth}
        \includegraphics[width=\linewidth]{nr_RIO_0.092.png}
        \caption{}
        \label{fig:rio_0.092_nr}
    \end{subfigure}%
    \begin{subfigure}{0.5\textwidth}
        \includegraphics[width=\linewidth]{gamma_tau_RIO_0.092.png}
        \caption{}
        \label{fig:rio_0.092_gamma_tau}
    \end{subfigure}
    \caption{Here are represented the important features about the phase transition in Rio De Janeiro for UCI = 0.092.
    In \ref{fig:rio_0.092_nt} is represented the number of people trapped in the road network at time t. In \ref{fig:rio_0.092_scaled_nt}, the vertical line representing $t_{\infty}$. we see the values $n(t)t^{\alpha}$ over $\epsilon (t-t_0)$, with $\epsilon = \frac{R^2 - R_c^2}{R_c^2}$ with t up to $t_{start} + t_{end_{fit}}$ In \ref{fig:rio_0.092_nr} showing the phase transition requirement. We plot n(R,$t_{\infty}$) with $\beta = $ 1.430. In \ref{fig:rio_0.092_gamma_tau} we see the linear increase of $\tau$ compared to $\Gamma$}
    \label{fig:phase_transition_rio_0.092}
\end{figure}

\begin{figure}
    \centering
    \begin{subfigure}{0.5\textwidth}
        \includegraphics[width=\linewidth]{nt_RIO_0.099.png}
        \caption{}
        \label{fig:rio_0.099_nt}
    \end{subfigure}%
    \begin{subfigure}{0.5\textwidth}
        \includegraphics[width=\linewidth]{scaled_nt_RIO_0.099.png}
        \caption{}
        \label{fig:rio_0.099_scaled_nt}
    \end{subfigure}
    
    \begin{subfigure}{0.5\textwidth}
        \includegraphics[width=\linewidth]{nr_RIO_0.099.png}
        \caption{}
        \label{fig:rio_0.099_nr}
    \end{subfigure}%
    \begin{subfigure}{0.5\textwidth}
        \includegraphics[width=\linewidth]{gamma_tau_RIO_0.099.png}
        \caption{}
        \label{fig:rio_0.099_gamma_tau}
    \end{subfigure}
    \caption{Here are represented the important features about the phase transition in Rio De Janeiro for UCI = 0.099.
    In \ref{fig:rio_0.099_nt} is represented the number of people trapped in the road network at time t. In \ref{fig:rio_0.099_scaled_nt}, the vertical line representing $t_{\infty}$. we see the values $n(t)t^{\alpha}$ over $\epsilon (t-t_0)$, with $\epsilon = \frac{R^2 - R_c^2}{R_c^2}$ with t up to $t_{start} + t_{end_{fit}}$ In \ref{fig:rio_0.099_nr} showing the phase transition requirement. We plot n(R,$t_{\infty}$) with $\beta = $ 0.055. In \ref{fig:rio_0.099_gamma_tau} we see the linear increase of $\tau$ compared to $\Gamma$}
    \label{fig:phase_transition_rio_0.099}
\end{figure}

\begin{figure}
    \centering
    \begin{subfigure}{0.5\textwidth}
        \includegraphics[width=\linewidth]{nt_RIO_0.102.png}
        \caption{}
        \label{fig:rio_0.102_nt}
    \end{subfigure}%
    \begin{subfigure}{0.5\textwidth}
        \includegraphics[width=\linewidth]{scaled_nt_RIO_0.102.png}
        \caption{}
        \label{fig:rio_0.102_scaled_nt}
    \end{subfigure}
    
    \begin{subfigure}{0.5\textwidth}
        \includegraphics[width=\linewidth]{nr_RIO_0.102.png}
        \caption{}
        \label{fig:rio_0.102_nr}
    \end{subfigure}%
    \begin{subfigure}{0.5\textwidth}
        \includegraphics[width=\linewidth]{gamma_tau_RIO_0.102.png}
        \caption{}
        \label{fig:rio_0.102_gamma_tau}
    \end{subfigure}
    \caption{Here are represented the important features about the phase transition in Rio De Janeiro for UCI = 0.102.
    In \ref{fig:rio_0.102_nt} is represented the number of people trapped in the road network at time t. In \ref{fig:rio_0.102_scaled_nt}, the vertical line representing $t_{\infty}$. we see the values $n(t)t^{\alpha}$ over $\epsilon (t-t_0)$, with $\epsilon = \frac{R^2 - R_c^2}{R_c^2}$ with t up to $t_{start} + t_{end_{fit}}$ In \ref{fig:rio_0.102_nr} showing the phase transition requirement. We plot n(R,$t_{\infty}$) with $\beta = $ 1.321. In \ref{fig:rio_0.102_gamma_tau} we see the linear increase of $\tau$ compared to $\Gamma$}
    \label{fig:phase_transition_rio_0.102}
\end{figure}

\begin{figure}
    \centering
    \begin{subfigure}{0.5\textwidth}
        \includegraphics[width=\linewidth]{nt_RIO_0.125.png}
        \caption{}
        \label{fig:rio_0.125_nt}
    \end{subfigure}%
    \begin{subfigure}{0.5\textwidth}
        \includegraphics[width=\linewidth]{scaled_nt_RIO_0.125.png}
        \caption{}
        \label{fig:rio_0.125_scaled_nt}
    \end{subfigure}
    
    \begin{subfigure}{0.5\textwidth}
        \includegraphics[width=\linewidth]{nr_RIO_0.125.png}
        \caption{}
        \label{fig:rio_0.125_nr}
    \end{subfigure}%
    \begin{subfigure}{0.5\textwidth}
        \includegraphics[width=\linewidth]{gamma_tau_RIO_0.125.png}
        \caption{}
        \label{fig:rio_0.125_gamma_tau}
    \end{subfigure}
    \caption{Here are represented the important features about the phase transition in Rio De Janeiro for UCI = 0.125.
    In \ref{fig:rio_0.125_nt} is represented the number of people trapped in the road network at time t. In \ref{fig:rio_0.125_scaled_nt}, the vertical line representing $t_{\infty}$. we see the values $n(t)t^{\alpha}$ over $\epsilon (t-t_0)$, with $\epsilon = \frac{R^2 - R_c^2}{R_c^2}$ with t up to $t_{start} + t_{end_{fit}}$ In \ref{fig:rio_0.125_nr} showing the phase transition requirement. We plot n(R,$t_{\infty}$) with $\beta = $ 0.427. In \ref{fig:rio_0.125_gamma_tau} we see the linear increase of $\tau$ compared to $\Gamma$}
    \label{fig:phase_transition_rio_0.125}
\end{figure}

\begin{figure}
    \centering
    \begin{subfigure}{0.5\textwidth}
        \includegraphics[width=\linewidth]{nt_RIO_0.127.png}
        \caption{}
        \label{fig:rio_0.127_nt}
    \end{subfigure}%
    \begin{subfigure}{0.5\textwidth}
        \includegraphics[width=\linewidth]{scaled_nt_RIO_0.127.png}
        \caption{}
        \label{fig:rio_0.127_scaled_nt}
    \end{subfigure}
    
    \begin{subfigure}{0.5\textwidth}
        \includegraphics[width=\linewidth]{nr_RIO_0.127.png}
        \caption{}
        \label{fig:rio_0.127_nr}
    \end{subfigure}%
    \begin{subfigure}{0.5\textwidth}
        \includegraphics[width=\linewidth]{gamma_tau_RIO_0.127.png}
        \caption{}
        \label{fig:rio_0.127_gamma_tau}
    \end{subfigure}
    \caption{Here are represented the important features about the phase transition in Rio De Janeiro for UCI = 0.127.
    In \ref{fig:rio_0.127_nt} is represented the number of people trapped in the road network at time t. In \ref{fig:rio_0.127_scaled_nt}, the vertical line representing $t_{\infty}$. we see the values $n(t)t^{\alpha}$ over $\epsilon (t-t_0)$, with $\epsilon = \frac{R^2 - R_c^2}{R_c^2}$ with t up to $t_{start} + t_{end_{fit}}$ In \ref{fig:rio_0.127_nr} showing the phase transition requirement. We plot n(R,$t_{\infty}$) with $\beta = $ 0.326. In \ref{fig:rio_0.127_gamma_tau} we see the linear increase of $\tau$ compared to $\Gamma$}
    \label{fig:phase_transition_rio_0.127}
\end{figure}

\begin{figure}
    \centering
    \begin{subfigure}{0.5\textwidth}
        \includegraphics[width=\linewidth]{nt_RIO_0.215.png}
        \caption{}
        \label{fig:rio_0.215_nt}
    \end{subfigure}%
    \begin{subfigure}{0.5\textwidth}
        \includegraphics[width=\linewidth]{scaled_nt_RIO_0.215.png}
        \caption{}
        \label{fig:rio_0.215_scaled_nt}
    \end{subfigure}
    
    \begin{subfigure}{0.5\textwidth}
        \includegraphics[width=\linewidth]{nr_RIO_0.215.png}
        \caption{}
        \label{fig:rio_0.215_nr}
    \end{subfigure}%
    \begin{subfigure}{0.5\textwidth}
        \includegraphics[width=\linewidth]{gamma_tau_RIO_0.215.png}
        \caption{}
        \label{fig:rio_0.215_gamma_tau}
    \end{subfigure}
    \caption{Here are represented the important features about the phase transition in Rio De Janeiro for UCI = 0.215.
    In \ref{fig:rio_0.215_nt} is represented the number of people trapped in the road network at time t. In \ref{fig:rio_0.215_scaled_nt}, the vertical line representing $t_{\infty}$. we see the values $n(t)t^{\alpha}$ over $\epsilon (t-t_0)$, with $\epsilon = \frac{R^2 - R_c^2}{R_c^2}$ with t up to $t_{start} + t_{end_{fit}}$ In \ref{fig:rio_0.215_nr} showing the phase transition requirement. We plot n(R,$t_{\infty}$) with $\beta = $ 0.035. In \ref{fig:rio_0.215_gamma_tau} we see the linear increase of $\tau$ compared to $\Gamma$}
    \label{fig:phase_transition_rio_0.215}
\end{figure}

\begin{figure}
    \centering
    \begin{subfigure}{0.5\textwidth}
        \includegraphics[width=\linewidth]{nt_RIO_0.241.png}
        \caption{}
        \label{fig:rio_0.241_nt}
    \end{subfigure}%
    \begin{subfigure}{0.5\textwidth}
        \includegraphics[width=\linewidth]{scaled_nt_RIO_0.241.png}
        \caption{}
        \label{fig:rio_0.241_scaled_nt}
    \end{subfigure}
    
    \begin{subfigure}{0.5\textwidth}
        \includegraphics[width=\linewidth]{nr_RIO_0.241.png}
        \caption{}
        \label{fig:rio_0.241_nr}
    \end{subfigure}%
    \begin{subfigure}{0.5\textwidth}
        \includegraphics[width=\linewidth]{gamma_tau_RIO_0.241.png}
        \caption{}
        \label{fig:rio_0.241_gamma_tau}
    \end{subfigure}
    \caption{Here are represented the important features about the phase transition in Rio De Janeiro for UCI = 0.241.
    In \ref{fig:rio_0.241_nt} is represented the number of people trapped in the road network at time t. In \ref{fig:rio_0.241_scaled_nt}, the vertical line representing $t_{\infty}$. we see the values $n(t)t^{\alpha}$ over $\epsilon (t-t_0)$, with $\epsilon = \frac{R^2 - R_c^2}{R_c^2}$ with t up to $t_{start} + t_{end_{fit}}$ In \ref{fig:rio_0.241_nr} showing the phase transition requirement. We plot n(R,$t_{\infty}$) with $\beta = $ 0.041. In \ref{fig:rio_0.241_gamma_tau} we see the linear increase of $\tau$ compared to $\Gamma$}
    \label{fig:phase_transition_rio_0.241}
\end{figure}

\begin{figure}
    \centering
    \begin{subfigure}{0.5\textwidth}
        \includegraphics[width=\linewidth]{nt_RIO_0.297.png}
        \caption{}
        \label{fig:rio_0.297_nt}
    \end{subfigure}%
    \begin{subfigure}{0.5\textwidth}
        \includegraphics[width=\linewidth]{scaled_nt_RIO_0.297.png}
        \caption{}
        \label{fig:rio_0.297_scaled_nt}
    \end{subfigure}
    
    \begin{subfigure}{0.5\textwidth}
        \includegraphics[width=\linewidth]{nr_RIO_0.297.png}
        \caption{}
        \label{fig:rio_0.297_nr}
    \end{subfigure}%
    \begin{subfigure}{0.5\textwidth}
        \includegraphics[width=\linewidth]{gamma_tau_RIO_0.297.png}
        \caption{}
        \label{fig:rio_0.297_gamma_tau}
    \end{subfigure}
    \caption{Here are represented the important features about the phase transition in Rio De Janeiro for UCI = 0.297.
    In \ref{fig:rio_0.297_nt} is represented the number of people trapped in the road network at time t. In \ref{fig:rio_0.297_scaled_nt}, the vertical line representing $t_{\infty}$. we see the values $n(t)t^{\alpha}$ over $\epsilon (t-t_0)$, with $\epsilon = \frac{R^2 - R_c^2}{R_c^2}$ with t up to $t_{start} + t_{end_{fit}}$ In \ref{fig:rio_0.297_nr} showing the phase transition requirement. We plot n(R,$t_{\infty}$) with $\beta = $ 0.091. In \ref{fig:rio_0.297_gamma_tau} we see the linear increase of $\tau$ compared to $\Gamma$}
    \label{fig:phase_transition_rio_0.297}
\end{figure}

\begin{figure}
    \centering
    \begin{subfigure}{0.5\textwidth}
        \includegraphics[width=\linewidth]{nt_RIO_0.298.png}
        \caption{}
        \label{fig:rio_0.298_nt}
    \end{subfigure}%
    \begin{subfigure}{0.5\textwidth}
        \includegraphics[width=\linewidth]{scaled_nt_RIO_0.298.png}
        \caption{}
        \label{fig:rio_0.298_scaled_nt}
    \end{subfigure}
    
    \begin{subfigure}{0.5\textwidth}
        \includegraphics[width=\linewidth]{nr_RIO_0.298.png}
        \caption{}
        \label{fig:rio_0.298_nr}
    \end{subfigure}%
    \begin{subfigure}{0.5\textwidth}
        \includegraphics[width=\linewidth]{gamma_tau_RIO_0.298.png}
        \caption{}
        \label{fig:rio_0.298_gamma_tau}
    \end{subfigure}
    \caption{Here are represented the important features about the phase transition in Rio De Janeiro for UCI = 0.298.
    In \ref{fig:rio_0.298_nt} is represented the number of people trapped in the road network at time t. In \ref{fig:rio_0.298_scaled_nt}, the vertical line representing $t_{\infty}$. we see the values $n(t)t^{\alpha}$ over $\epsilon (t-t_0)$, with $\epsilon = \frac{R^2 - R_c^2}{R_c^2}$ with t up to $t_{start} + t_{end_{fit}}$ In \ref{fig:rio_0.298_nr} showing the phase transition requirement. We plot n(R,$t_{\infty}$) with $\beta = $ 0.134. In \ref{fig:rio_0.298_gamma_tau} we see the linear increase of $\tau$ compared to $\Gamma$}
    \label{fig:phase_transition_rio_0.298}
\end{figure}

\begin{figure}
    \centering
    \begin{subfigure}{0.5\textwidth}
        \includegraphics[width=\linewidth]{nt_RIO_0.315.png}
        \caption{}
        \label{fig:rio_0.315_nt}
    \end{subfigure}%
    \begin{subfigure}{0.5\textwidth}
        \includegraphics[width=\linewidth]{scaled_nt_RIO_0.315.png}
        \caption{}
        \label{fig:rio_0.315_scaled_nt}
    \end{subfigure}
    
    \begin{subfigure}{0.5\textwidth}
        \includegraphics[width=\linewidth]{nr_RIO_0.315.png}
        \caption{}
        \label{fig:rio_0.315_nr}
    \end{subfigure}%
    \begin{subfigure}{0.5\textwidth}
        \includegraphics[width=\linewidth]{gamma_tau_RIO_0.315.png}
        \caption{}
        \label{fig:rio_0.315_gamma_tau}
    \end{subfigure}
    \caption{Here are represented the important features about the phase transition in Rio De Janeiro for UCI = 0.315.
    In \ref{fig:rio_0.315_nt} is represented the number of people trapped in the road network at time t. In \ref{fig:rio_0.315_scaled_nt}, the vertical line representing $t_{\infty}$. we see the values $n(t)t^{\alpha}$ over $\epsilon (t-t_0)$, with $\epsilon = \frac{R^2 - R_c^2}{R_c^2}$ with t up to $t_{start} + t_{end_{fit}}$ In \ref{fig:rio_0.315_nr} showing the phase transition requirement. We plot n(R,$t_{\infty}$) with $\beta = $ 0.071. In \ref{fig:rio_0.315_gamma_tau} we see the linear increase of $\tau$ compared to $\Gamma$}
    \label{fig:phase_transition_rio_0.315}
\end{figure}

\begin{figure}
    \centering
    \begin{subfigure}{0.5\textwidth}
        \includegraphics[width=\linewidth]{nt_RIO_0.317.png}
        \caption{}
        \label{fig:rio_0.317_nt}
    \end{subfigure}%
    \begin{subfigure}{0.5\textwidth}
        \includegraphics[width=\linewidth]{scaled_nt_RIO_0.317.png}
        \caption{}
        \label{fig:rio_0.317_scaled_nt}
    \end{subfigure}
    
    \begin{subfigure}{0.5\textwidth}
        \includegraphics[width=\linewidth]{nr_RIO_0.317.png}
        \caption{}
        \label{fig:rio_0.317_nr}
    \end{subfigure}%
    \begin{subfigure}{0.5\textwidth}
        \includegraphics[width=\linewidth]{gamma_tau_RIO_0.317.png}
        \caption{}
        \label{fig:rio_0.317_gamma_tau}
    \end{subfigure}
    \caption{Here are represented the important features about the phase transition in Rio De Janeiro for UCI = 0.317.
    In \ref{fig:rio_0.317_nt} is represented the number of people trapped in the road network at time t. In \ref{fig:rio_0.317_scaled_nt}, the vertical line representing $t_{\infty}$. we see the values $n(t)t^{\alpha}$ over $\epsilon (t-t_0)$, with $\epsilon = \frac{R^2 - R_c^2}{R_c^2}$ with t up to $t_{start} + t_{end_{fit}}$ In \ref{fig:rio_0.317_nr} showing the phase transition requirement. We plot n(R,$t_{\infty}$) with $\beta = $ 0.043. In \ref{fig:rio_0.317_gamma_tau} we see the linear increase of $\tau$ compared to $\Gamma$}
    \label{fig:phase_transition_rio_0.317}
\end{figure}

\begin{figure}
    \centering
    \begin{subfigure}{0.5\textwidth}
        \includegraphics[width=\linewidth]{nt_RIO_0.348.png}
        \caption{}
        \label{fig:rio_0.348_nt}
    \end{subfigure}%
    \begin{subfigure}{0.5\textwidth}
        \includegraphics[width=\linewidth]{scaled_nt_RIO_0.348.png}
        \caption{}
        \label{fig:rio_0.348_scaled_nt}
    \end{subfigure}
    
    \begin{subfigure}{0.5\textwidth}
        \includegraphics[width=\linewidth]{nr_RIO_0.348.png}
        \caption{}
        \label{fig:rio_0.348_nr}
    \end{subfigure}%
    \begin{subfigure}{0.5\textwidth}
        \includegraphics[width=\linewidth]{gamma_tau_RIO_0.348.png}
        \caption{}
        \label{fig:rio_0.348_gamma_tau}
    \end{subfigure}
    \caption{Here are represented the important features about the phase transition in Rio De Janeiro for UCI = 0.348.
    In \ref{fig:rio_0.348_nt} is represented the number of people trapped in the road network at time t. In \ref{fig:rio_0.348_scaled_nt}, the vertical line representing $t_{\infty}$. we see the values $n(t)t^{\alpha}$ over $\epsilon (t-t_0)$, with $\epsilon = \frac{R^2 - R_c^2}{R_c^2}$ with t up to $t_{start} + t_{end_{fit}}$ In \ref{fig:rio_0.348_nr} showing the phase transition requirement. We plot n(R,$t_{\infty}$) with $\beta = $ 0.053. In \ref{fig:rio_0.348_gamma_tau} we see the linear increase of $\tau$ compared to $\Gamma$}
    \label{fig:phase_transition_rio_0.348}
\end{figure}

\begin{figure}
    \centering
    \begin{subfigure}{0.5\textwidth}
        \includegraphics[width=\linewidth]{nt_RIO_0.349.png}
        \caption{}
        \label{fig:rio_0.349_nt}
    \end{subfigure}%
    \begin{subfigure}{0.5\textwidth}
        \includegraphics[width=\linewidth]{scaled_nt_RIO_0.349.png}
        \caption{}
        \label{fig:rio_0.349_scaled_nt}
    \end{subfigure}
    
    \begin{subfigure}{0.5\textwidth}
        \includegraphics[width=\linewidth]{nr_RIO_0.349.png}
        \caption{}
        \label{fig:rio_0.349_nr}
    \end{subfigure}%
    \begin{subfigure}{0.5\textwidth}
        \includegraphics[width=\linewidth]{gamma_tau_RIO_0.349.png}
        \caption{}
        \label{fig:rio_0.349_gamma_tau}
    \end{subfigure}
    \caption{Here are represented the important features about the phase transition in Rio De Janeiro for UCI = 0.349.
    In \ref{fig:rio_0.349_nt} is represented the number of people trapped in the road network at time t. In \ref{fig:rio_0.349_scaled_nt}, the vertical line representing $t_{\infty}$. we see the values $n(t)t^{\alpha}$ over $\epsilon (t-t_0)$, with $\epsilon = \frac{R^2 - R_c^2}{R_c^2}$ with t up to $t_{start} + t_{end_{fit}}$ In \ref{fig:rio_0.349_nr} showing the phase transition requirement. We plot n(R,$t_{\infty}$) with $\beta = $ 0.044. In \ref{fig:rio_0.349_gamma_tau} we see the linear increase of $\tau$ compared to $\Gamma$}
    \label{fig:phase_transition_rio_0.349}
\end{figure}

\begin{figure}
    \centering
    \begin{subfigure}{0.5\textwidth}
        \includegraphics[width=\linewidth]{nt_RIO_0.413.png}
        \caption{}
        \label{fig:rio_0.413_nt}
    \end{subfigure}%
    \begin{subfigure}{0.5\textwidth}
        \includegraphics[width=\linewidth]{scaled_nt_RIO_0.413.png}
        \caption{}
        \label{fig:rio_0.413_scaled_nt}
    \end{subfigure}
    
    \begin{subfigure}{0.5\textwidth}
        \includegraphics[width=\linewidth]{nr_RIO_0.413.png}
        \caption{}
        \label{fig:rio_0.413_nr}
    \end{subfigure}%
    \begin{subfigure}{0.5\textwidth}
        \includegraphics[width=\linewidth]{gamma_tau_RIO_0.413.png}
        \caption{}
        \label{fig:rio_0.413_gamma_tau}
    \end{subfigure}
    \caption{Here are represented the important features about the phase transition in Rio De Janeiro for UCI = 0.413.
    In \ref{fig:rio_0.413_nt} is represented the number of people trapped in the road network at time t. In \ref{fig:rio_0.413_scaled_nt}, the vertical line representing $t_{\infty}$. we see the values $n(t)t^{\alpha}$ over $\epsilon (t-t_0)$, with $\epsilon = \frac{R^2 - R_c^2}{R_c^2}$ with t up to $t_{start} + t_{end_{fit}}$ In \ref{fig:rio_0.413_nr} showing the phase transition requirement. We plot n(R,$t_{\infty}$) with $\beta = $ 0.092. In \ref{fig:rio_0.413_gamma_tau} we see the linear increase of $\tau$ compared to $\Gamma$}
    \label{fig:phase_transition_rio_0.413}
\end{figure}

\begin{figure}
    \centering
    \begin{subfigure}{0.5\textwidth}
        \includegraphics[width=\linewidth]{nt_RIO_0.458.png}
        \caption{}
        \label{fig:rio_0.458_nt}
    \end{subfigure}%
    \begin{subfigure}{0.5\textwidth}
        \includegraphics[width=\linewidth]{scaled_nt_RIO_0.458.png}
        \caption{}
        \label{fig:rio_0.458_scaled_nt}
    \end{subfigure}
    
    \begin{subfigure}{0.5\textwidth}
        \includegraphics[width=\linewidth]{nr_RIO_0.458.png}
        \caption{}
        \label{fig:rio_0.458_nr}
    \end{subfigure}%
    \begin{subfigure}{0.5\textwidth}
        \includegraphics[width=\linewidth]{gamma_tau_RIO_0.458.png}
        \caption{}
        \label{fig:rio_0.458_gamma_tau}
    \end{subfigure}
    \caption{Here are represented the important features about the phase transition in Rio De Janeiro for UCI = 0.458.
    In \ref{fig:rio_0.458_nt} is represented the number of people trapped in the road network at time t. In \ref{fig:rio_0.458_scaled_nt}, the vertical line representing $t_{\infty}$. we see the values $n(t)t^{\alpha}$ over $\epsilon (t-t_0)$, with $\epsilon = \frac{R^2 - R_c^2}{R_c^2}$ with t up to $t_{start} + t_{end_{fit}}$ In \ref{fig:rio_0.458_nr} showing the phase transition requirement. We plot n(R,$t_{\infty}$) with $\beta = $ 0.067. In \ref{fig:rio_0.458_gamma_tau} we see the linear increase of $\tau$ compared to $\Gamma$}
    \label{fig:phase_transition_rio_0.458}
\end{figure}

\begin{figure}
    \centering
    \begin{subfigure}{0.5\textwidth}
        \includegraphics[width=\linewidth]{nt_RIO_0.543.png}
        \caption{}
        \label{fig:rio_0.543_nt}
    \end{subfigure}%
    \begin{subfigure}{0.5\textwidth}
        \includegraphics[width=\linewidth]{scaled_nt_RIO_0.543.png}
        \caption{}
        \label{fig:rio_0.543_scaled_nt}
    \end{subfigure}
    
    \begin{subfigure}{0.5\textwidth}
        \includegraphics[width=\linewidth]{nr_RIO_0.543.png}
        \caption{}
        \label{fig:rio_0.543_nr}
    \end{subfigure}%
    \begin{subfigure}{0.5\textwidth}
        \includegraphics[width=\linewidth]{gamma_tau_RIO_0.543.png}
        \caption{}
        \label{fig:rio_0.543_gamma_tau}
    \end{subfigure}
    \caption{Here are represented the important features about the phase transition in Rio De Janeiro for UCI = 0.543.
    In \ref{fig:rio_0.543_nt} is represented the number of people trapped in the road network at time t. In \ref{fig:rio_0.543_scaled_nt}, the vertical line representing $t_{\infty}$. we see the values $n(t)t^{\alpha}$ over $\epsilon (t-t_0)$, with $\epsilon = \frac{R^2 - R_c^2}{R_c^2}$ with t up to $t_{start} + t_{end_{fit}}$ In \ref{fig:rio_0.543_nr} showing the phase transition requirement. We plot n(R,$t_{\infty}$) with $\beta = $ 0.102. In \ref{fig:rio_0.543_gamma_tau} we see the linear increase of $\tau$ compared to $\Gamma$}
    \label{fig:phase_transition_rio_0.543}
\end{figure}

\begin{figure}
    \centering
    \begin{subfigure}{0.5\textwidth}
        \includegraphics[width=\linewidth]{nt_RIO_0.561.png}
        \caption{}
        \label{fig:rio_0.561_nt}
    \end{subfigure}%
    \begin{subfigure}{0.5\textwidth}
        \includegraphics[width=\linewidth]{scaled_nt_RIO_0.561.png}
        \caption{}
        \label{fig:rio_0.561_scaled_nt}
    \end{subfigure}
    
    \begin{subfigure}{0.5\textwidth}
        \includegraphics[width=\linewidth]{nr_RIO_0.561.png}
        \caption{}
        \label{fig:rio_0.561_nr}
    \end{subfigure}%
    \begin{subfigure}{0.5\textwidth}
        \includegraphics[width=\linewidth]{gamma_tau_RIO_0.561.png}
        \caption{}
        \label{fig:rio_0.561_gamma_tau}
    \end{subfigure}
    \caption{Here are represented the important features about the phase transition in Rio De Janeiro for UCI = 0.561.
    In \ref{fig:rio_0.561_nt} is represented the number of people trapped in the road network at time t. In \ref{fig:rio_0.561_scaled_nt}, the vertical line representing $t_{\infty}$. we see the values $n(t)t^{\alpha}$ over $\epsilon (t-t_0)$, with $\epsilon = \frac{R^2 - R_c^2}{R_c^2}$ with t up to $t_{start} + t_{end_{fit}}$ In \ref{fig:rio_0.561_nr} showing the phase transition requirement. We plot n(R,$t_{\infty}$) with $\beta = $ 0.032. In \ref{fig:rio_0.561_gamma_tau} we see the linear increase of $\tau$ compared to $\Gamma$}
    \label{fig:phase_transition_rio_0.561}
\end{figure}

\begin{figure}
    \centering
    \begin{subfigure}{0.5\textwidth}
        \includegraphics[width=\linewidth]{nt_RIO_0.581.png}
        \caption{}
        \label{fig:rio_0.581_nt}
    \end{subfigure}%
    \begin{subfigure}{0.5\textwidth}
        \includegraphics[width=\linewidth]{scaled_nt_RIO_0.581.png}
        \caption{}
        \label{fig:rio_0.581_scaled_nt}
    \end{subfigure}
    
    \begin{subfigure}{0.5\textwidth}
        \includegraphics[width=\linewidth]{nr_RIO_0.581.png}
        \caption{}
        \label{fig:rio_0.581_nr}
    \end{subfigure}%
    \begin{subfigure}{0.5\textwidth}
        \includegraphics[width=\linewidth]{gamma_tau_RIO_0.581.png}
        \caption{}
        \label{fig:rio_0.581_gamma_tau}
    \end{subfigure}
    \caption{Here are represented the important features about the phase transition in Rio De Janeiro for UCI = 0.581.
    In \ref{fig:rio_0.581_nt} is represented the number of people trapped in the road network at time t. In \ref{fig:rio_0.581_scaled_nt}, the vertical line representing $t_{\infty}$. we see the values $n(t)t^{\alpha}$ over $\epsilon (t-t_0)$, with $\epsilon = \frac{R^2 - R_c^2}{R_c^2}$ with t up to $t_{start} + t_{end_{fit}}$ In \ref{fig:rio_0.581_nr} showing the phase transition requirement. We plot n(R,$t_{\infty}$) with $\beta = $ 0.281. In \ref{fig:rio_0.581_gamma_tau} we see the linear increase of $\tau$ compared to $\Gamma$}
    \label{fig:phase_transition_rio_0.581}
\end{figure}

\begin{figure}
    \centering
    \begin{subfigure}{0.5\textwidth}
        \includegraphics[width=\linewidth]{nt_RIO_0.679.png}
        \caption{}
        \label{fig:rio_0.679_nt}
    \end{subfigure}%
    \begin{subfigure}{0.5\textwidth}
        \includegraphics[width=\linewidth]{scaled_nt_RIO_0.679.png}
        \caption{}
        \label{fig:rio_0.679_scaled_nt}
    \end{subfigure}
    
    \begin{subfigure}{0.5\textwidth}
        \includegraphics[width=\linewidth]{nr_RIO_0.679.png}
        \caption{}
        \label{fig:rio_0.679_nr}
    \end{subfigure}%
    \begin{subfigure}{0.5\textwidth}
        \includegraphics[width=\linewidth]{gamma_tau_RIO_0.679.png}
        \caption{}
        \label{fig:rio_0.679_gamma_tau}
    \end{subfigure}
    \caption{Here are represented the important features about the phase transition in Rio De Janeiro for UCI = 0.679.
    In \ref{fig:rio_0.679_nt} is represented the number of people trapped in the road network at time t. In \ref{fig:rio_0.679_scaled_nt}, the vertical line representing $t_{\infty}$. we see the values $n(t)t^{\alpha}$ over $\epsilon (t-t_0)$, with $\epsilon = \frac{R^2 - R_c^2}{R_c^2}$ with t up to $t_{start} + t_{end_{fit}}$ In \ref{fig:rio_0.679_nr} showing the phase transition requirement. We plot n(R,$t_{\infty}$) with $\beta = $ 0.289. In \ref{fig:rio_0.679_gamma_tau} we see the linear increase of $\tau$ compared to $\Gamma$}
    \label{fig:phase_transition_rio_0.679}
\end{figure}

\begin{figure}
    \centering
    \begin{subfigure}{0.5\textwidth}
        \includegraphics[width=\linewidth]{nt_RIO_0.683.png}
        \caption{}
        \label{fig:rio_0.683_nt}
    \end{subfigure}%
    \begin{subfigure}{0.5\textwidth}
        \includegraphics[width=\linewidth]{scaled_nt_RIO_0.683.png}
        \caption{}
        \label{fig:rio_0.683_scaled_nt}
    \end{subfigure}
    
    \begin{subfigure}{0.5\textwidth}
        \includegraphics[width=\linewidth]{nr_RIO_0.683.png}
        \caption{}
        \label{fig:rio_0.683_nr}
    \end{subfigure}%
    \begin{subfigure}{0.5\textwidth}
        \includegraphics[width=\linewidth]{gamma_tau_RIO_0.683.png}
        \caption{}
        \label{fig:rio_0.683_gamma_tau}
    \end{subfigure}
    \caption{Here are represented the important features about the phase transition in Rio De Janeiro for UCI = 0.683.
    In \ref{fig:rio_0.683_nt} is represented the number of people trapped in the road network at time t. In \ref{fig:rio_0.683_scaled_nt}, the vertical line representing $t_{\infty}$. we see the values $n(t)t^{\alpha}$ over $\epsilon (t-t_0)$, with $\epsilon = \frac{R^2 - R_c^2}{R_c^2}$ with t up to $t_{start} + t_{end_{fit}}$ In \ref{fig:rio_0.683_nr} showing the phase transition requirement. We plot n(R,$t_{\infty}$) with $\beta = $ 1.107. In \ref{fig:rio_0.683_gamma_tau} we see the linear increase of $\tau$ compared to $\Gamma$}
    \label{fig:phase_transition_rio_0.683}
\end{figure}

\begin{figure}
    \centering
    \begin{subfigure}{0.5\textwidth}
        \includegraphics[width=\linewidth]{nt_RIO_0.693.png}
        \caption{}
        \label{fig:rio_0.693_nt}
    \end{subfigure}%
    \begin{subfigure}{0.5\textwidth}
        \includegraphics[width=\linewidth]{scaled_nt_RIO_0.693.png}
        \caption{}
        \label{fig:rio_0.693_scaled_nt}
    \end{subfigure}
    
    \begin{subfigure}{0.5\textwidth}
        \includegraphics[width=\linewidth]{nr_RIO_0.693.png}
        \caption{}
        \label{fig:rio_0.693_nr}
    \end{subfigure}%
    \begin{subfigure}{0.5\textwidth}
        \includegraphics[width=\linewidth]{gamma_tau_RIO_0.693.png}
        \caption{}
        \label{fig:rio_0.693_gamma_tau}
    \end{subfigure}
    \caption{Here are represented the important features about the phase transition in Rio De Janeiro for UCI = 0.693.
    In \ref{fig:rio_0.693_nt} is represented the number of people trapped in the road network at time t. In \ref{fig:rio_0.693_scaled_nt}, the vertical line representing $t_{\infty}$. we see the values $n(t)t^{\alpha}$ over $\epsilon (t-t_0)$, with $\epsilon = \frac{R^2 - R_c^2}{R_c^2}$ with t up to $t_{start} + t_{end_{fit}}$ In \ref{fig:rio_0.693_nr} showing the phase transition requirement. We plot n(R,$t_{\infty}$) with $\beta = $ 0.069. In \ref{fig:rio_0.693_gamma_tau} we see the linear increase of $\tau$ compared to $\Gamma$}
    \label{fig:phase_transition_rio_0.693}
\end{figure}

\begin{figure}
    \centering
    \begin{subfigure}{0.5\textwidth}
        \includegraphics[width=\linewidth]{nt_RIO_0.779.png}
        \caption{}
        \label{fig:rio_0.779_nt}
    \end{subfigure}%
    \begin{subfigure}{0.5\textwidth}
        \includegraphics[width=\linewidth]{scaled_nt_RIO_0.779.png}
        \caption{}
        \label{fig:rio_0.779_scaled_nt}
    \end{subfigure}
    
    \begin{subfigure}{0.5\textwidth}
        \includegraphics[width=\linewidth]{nr_RIO_0.779.png}
        \caption{}
        \label{fig:rio_0.779_nr}
    \end{subfigure}%
    \begin{subfigure}{0.5\textwidth}
        \includegraphics[width=\linewidth]{gamma_tau_RIO_0.779.png}
        \caption{}
        \label{fig:rio_0.779_gamma_tau}
    \end{subfigure}
    \caption{Here are represented the important features about the phase transition in Rio De Janeiro for UCI = 0.779.
    In \ref{fig:rio_0.779_nt} is represented the number of people trapped in the road network at time t. In \ref{fig:rio_0.779_scaled_nt}, the vertical line representing $t_{\infty}$. we see the values $n(t)t^{\alpha}$ over $\epsilon (t-t_0)$, with $\epsilon = \frac{R^2 - R_c^2}{R_c^2}$ with t up to $t_{start} + t_{end_{fit}}$ In \ref{fig:rio_0.779_nr} showing the phase transition requirement. We plot n(R,$t_{\infty}$) with $\beta = $ 0.048. In \ref{fig:rio_0.779_gamma_tau} we see the linear increase of $\tau$ compared to $\Gamma$}
    \label{fig:phase_transition_rio_0.779}
\end{figure}

\begin{figure}
    \centering
    \begin{subfigure}{0.5\textwidth}
        \includegraphics[width=\linewidth]{nt_RIO_0.789.png}
        \caption{}
        \label{fig:rio_0.789_nt}
    \end{subfigure}%
    \begin{subfigure}{0.5\textwidth}
        \includegraphics[width=\linewidth]{scaled_nt_RIO_0.789.png}
        \caption{}
        \label{fig:rio_0.789_scaled_nt}
    \end{subfigure}
    
    \begin{subfigure}{0.5\textwidth}
        \includegraphics[width=\linewidth]{nr_RIO_0.789.png}
        \caption{}
        \label{fig:rio_0.789_nr}
    \end{subfigure}%
    \begin{subfigure}{0.5\textwidth}
        \includegraphics[width=\linewidth]{gamma_tau_RIO_0.789.png}
        \caption{}
        \label{fig:rio_0.789_gamma_tau}
    \end{subfigure}
    \caption{Here are represented the important features about the phase transition in Rio De Janeiro for UCI = 0.789.
    In \ref{fig:rio_0.789_nt} is represented the number of people trapped in the road network at time t. In \ref{fig:rio_0.789_scaled_nt}, the vertical line representing $t_{\infty}$. we see the values $n(t)t^{\alpha}$ over $\epsilon (t-t_0)$, with $\epsilon = \frac{R^2 - R_c^2}{R_c^2}$ with t up to $t_{start} + t_{end_{fit}}$ In \ref{fig:rio_0.789_nr} showing the phase transition requirement. We plot n(R,$t_{\infty}$) with $\beta = $ 0.289. In \ref{fig:rio_0.789_gamma_tau} we see the linear increase of $\tau$ compared to $\Gamma$}
    \label{fig:phase_transition_rio_0.789}
\end{figure}

\begin{figure}
    \centering
    \begin{subfigure}{0.5\textwidth}
        \includegraphics[width=\linewidth]{nt_RIO_0.79.png}
        \caption{}
        \label{fig:rio_0.79_nt}
    \end{subfigure}%
    \begin{subfigure}{0.5\textwidth}
        \includegraphics[width=\linewidth]{scaled_nt_RIO_0.79.png}
        \caption{}
        \label{fig:rio_0.79_scaled_nt}
    \end{subfigure}
    
    \begin{subfigure}{0.5\textwidth}
        \includegraphics[width=\linewidth]{nr_RIO_0.79.png}
        \caption{}
        \label{fig:rio_0.79_nr}
    \end{subfigure}%
    \begin{subfigure}{0.5\textwidth}
        \includegraphics[width=\linewidth]{gamma_tau_RIO_0.79.png}
        \caption{}
        \label{fig:rio_0.79_gamma_tau}
    \end{subfigure}
    \caption{Here are represented the important features about the phase transition in Rio De Janeiro for UCI = 0.79.
    In \ref{fig:rio_0.79_nt} is represented the number of people trapped in the road network at time t. In \ref{fig:rio_0.79_scaled_nt}, the vertical line representing $t_{\infty}$. we see the values $n(t)t^{\alpha}$ over $\epsilon (t-t_0)$, with $\epsilon = \frac{R^2 - R_c^2}{R_c^2}$ with t up to $t_{start} + t_{end_{fit}}$ In \ref{fig:rio_0.79_nr} showing the phase transition requirement. We plot n(R,$t_{\infty}$) with $\beta = $ 0.433. In \ref{fig:rio_0.79_gamma_tau} we see the linear increase of $\tau$ compared to $\Gamma$}
    \label{fig:phase_transition_rio_0.79}
\end{figure}

\begin{figure}
    \centering
    \begin{subfigure}{0.5\textwidth}
        \includegraphics[width=\linewidth]{nt_RIO_0.82.png}
        \caption{}
        \label{fig:rio_0.82_nt}
    \end{subfigure}%
    \begin{subfigure}{0.5\textwidth}
        \includegraphics[width=\linewidth]{scaled_nt_RIO_0.82.png}
        \caption{}
        \label{fig:rio_0.82_scaled_nt}
    \end{subfigure}
    
    \begin{subfigure}{0.5\textwidth}
        \includegraphics[width=\linewidth]{nr_RIO_0.82.png}
        \caption{}
        \label{fig:rio_0.82_nr}
    \end{subfigure}%
    \begin{subfigure}{0.5\textwidth}
        \includegraphics[width=\linewidth]{gamma_tau_RIO_0.82.png}
        \caption{}
        \label{fig:rio_0.82_gamma_tau}
    \end{subfigure}
    \caption{Here are represented the important features about the phase transition in Rio De Janeiro for UCI = 0.82.
    In \ref{fig:rio_0.82_nt} is represented the number of people trapped in the road network at time t. In \ref{fig:rio_0.82_scaled_nt}, the vertical line representing $t_{\infty}$. we see the values $n(t)t^{\alpha}$ over $\epsilon (t-t_0)$, with $\epsilon = \frac{R^2 - R_c^2}{R_c^2}$ with t up to $t_{start} + t_{end_{fit}}$ In \ref{fig:rio_0.82_nr} showing the phase transition requirement. We plot n(R,$t_{\infty}$) with $\beta = $ 0.104. In \ref{fig:rio_0.82_gamma_tau} we see the linear increase of $\tau$ compared to $\Gamma$}
    \label{fig:phase_transition_rio_0.82}
\end{figure}

\begin{figure}
    \centering
    \begin{subfigure}{0.5\textwidth}
        \includegraphics[width=\linewidth]{nt_RIO_0.88.png}
        \caption{}
        \label{fig:rio_0.88_nt}
    \end{subfigure}%
    \begin{subfigure}{0.5\textwidth}
        \includegraphics[width=\linewidth]{scaled_nt_RIO_0.88.png}
        \caption{}
        \label{fig:rio_0.88_scaled_nt}
    \end{subfigure}
    
    \begin{subfigure}{0.5\textwidth}
        \includegraphics[width=\linewidth]{nr_RIO_0.88.png}
        \caption{}
        \label{fig:rio_0.88_nr}
    \end{subfigure}%
    \begin{subfigure}{0.5\textwidth}
        \includegraphics[width=\linewidth]{gamma_tau_RIO_0.88.png}
        \caption{}
        \label{fig:rio_0.88_gamma_tau}
    \end{subfigure}
    \caption{Here are represented the important features about the phase transition in Rio De Janeiro for UCI = 0.88.
    In \ref{fig:rio_0.88_nt} is represented the number of people trapped in the road network at time t. In \ref{fig:rio_0.88_scaled_nt}, the vertical line representing $t_{\infty}$. we see the values $n(t)t^{\alpha}$ over $\epsilon (t-t_0)$, with $\epsilon = \frac{R^2 - R_c^2}{R_c^2}$ with t up to $t_{start} + t_{end_{fit}}$ In \ref{fig:rio_0.88_nr} showing the phase transition requirement. We plot n(R,$t_{\infty}$) with $\beta = $ 0.630. In \ref{fig:rio_0.88_gamma_tau} we see the linear increase of $\tau$ compared to $\Gamma$}
    \label{fig:phase_transition_rio_0.88}
\end{figure}

\begin{figure}
    \centering
    \begin{subfigure}{0.5\textwidth}
        \includegraphics[width=\linewidth]{nt_RIO_0.904.png}
        \caption{}
        \label{fig:rio_0.904_nt}
    \end{subfigure}%
    \begin{subfigure}{0.5\textwidth}
        \includegraphics[width=\linewidth]{scaled_nt_RIO_0.904.png}
        \caption{}
        \label{fig:rio_0.904_scaled_nt}
    \end{subfigure}
    
    \begin{subfigure}{0.5\textwidth}
        \includegraphics[width=\linewidth]{nr_RIO_0.904.png}
        \caption{}
        \label{fig:rio_0.904_nr}
    \end{subfigure}%
    \begin{subfigure}{0.5\textwidth}
        \includegraphics[width=\linewidth]{gamma_tau_RIO_0.904.png}
        \caption{}
        \label{fig:rio_0.904_gamma_tau}
    \end{subfigure}
    \caption{Here are represented the important features about the phase transition in Rio De Janeiro for UCI = 0.904.
    In \ref{fig:rio_0.904_nt} is represented the number of people trapped in the road network at time t. In \ref{fig:rio_0.904_scaled_nt}, the vertical line representing $t_{\infty}$. we see the values $n(t)t^{\alpha}$ over $\epsilon (t-t_0)$, with $\epsilon = \frac{R^2 - R_c^2}{R_c^2}$ with t up to $t_{start} + t_{end_{fit}}$ In \ref{fig:rio_0.904_nr} showing the phase transition requirement. We plot n(R,$t_{\infty}$) with $\beta = $ 0.242. In \ref{fig:rio_0.904_gamma_tau} we see the linear increase of $\tau$ compared to $\Gamma$}
    \label{fig:phase_transition_rio_0.904}
\end{figure}

\begin{figure}
    \centering
    \begin{subfigure}{0.5\textwidth}
        \includegraphics[width=\linewidth]{nt_RIO_0.979.png}
        \caption{}
        \label{fig:rio_0.979_nt}
    \end{subfigure}%
    \begin{subfigure}{0.5\textwidth}
        \includegraphics[width=\linewidth]{scaled_nt_RIO_0.979.png}
        \caption{}
        \label{fig:rio_0.979_scaled_nt}
    \end{subfigure}
    
    \begin{subfigure}{0.5\textwidth}
        \includegraphics[width=\linewidth]{nr_RIO_0.979.png}
        \caption{}
        \label{fig:rio_0.979_nr}
    \end{subfigure}%
    \begin{subfigure}{0.5\textwidth}
        \includegraphics[width=\linewidth]{gamma_tau_RIO_0.979.png}
        \caption{}
        \label{fig:rio_0.979_gamma_tau}
    \end{subfigure}
    \caption{Here are represented the important features about the phase transition in Rio De Janeiro for UCI = 0.979.
    In \ref{fig:rio_0.979_nt} is represented the number of people trapped in the road network at time t. In \ref{fig:rio_0.979_scaled_nt}, the vertical line representing $t_{\infty}$. we see the values $n(t)t^{\alpha}$ over $\epsilon (t-t_0)$, with $\epsilon = \frac{R^2 - R_c^2}{R_c^2}$ with t up to $t_{start} + t_{end_{fit}}$ In \ref{fig:rio_0.979_nr} showing the phase transition requirement. We plot n(R,$t_{\infty}$) with $\beta = $ 0.045. In \ref{fig:rio_0.979_gamma_tau} we see the linear increase of $\tau$ compared to $\Gamma$}
    \label{fig:phase_transition_rio_0.979}
\end{figure}

\begin{figure}
    \centering
    \begin{subfigure}{0.5\textwidth}
        \includegraphics[width=\linewidth]{nt_RIO_0.999.png}
        \caption{}
        \label{fig:rio_0.999_nt}
    \end{subfigure}%
    \begin{subfigure}{0.5\textwidth}
        \includegraphics[width=\linewidth]{scaled_nt_RIO_0.999.png}
        \caption{}
        \label{fig:rio_0.999_scaled_nt}
    \end{subfigure}
    
    \begin{subfigure}{0.5\textwidth}
        \includegraphics[width=\linewidth]{nr_RIO_0.999.png}
        \caption{}
        \label{fig:rio_0.999_nr}
    \end{subfigure}%
    \begin{subfigure}{0.5\textwidth}
        \includegraphics[width=\linewidth]{gamma_tau_RIO_0.999.png}
        \caption{}
        \label{fig:rio_0.999_gamma_tau}
    \end{subfigure}
    \caption{Here are represented the important features about the phase transition in Rio De Janeiro for UCI = 0.999.
    In \ref{fig:rio_0.999_nt} is represented the number of people trapped in the road network at time t. In \ref{fig:rio_0.999_scaled_nt}, the vertical line representing $t_{\infty}$. we see the values $n(t)t^{\alpha}$ over $\epsilon (t-t_0)$, with $\epsilon = \frac{R^2 - R_c^2}{R_c^2}$ with t up to $t_{start} + t_{end_{fit}}$ In \ref{fig:rio_0.999_nr} showing the phase transition requirement. We plot n(R,$t_{\infty}$) with $\beta = $ 0.285. In \ref{fig:rio_0.999_gamma_tau} we see the linear increase of $\tau$ compared to $\Gamma$}
    \label{fig:phase_transition_rio_0.999}
\end{figure}

\begin{figure}
    \centering
    \begin{subfigure}{0.5\textwidth}
        \includegraphics[width=\linewidth]{nt_SFO_0.445.png}
        \caption{}
        \label{fig:sfo_0.445_nt}
    \end{subfigure}%
    \begin{subfigure}{0.5\textwidth}
        \includegraphics[width=\linewidth]{scaled_nt_SFO_0.445.png}
        \caption{}
        \label{fig:sfo_0.445_scaled_nt}
    \end{subfigure}
    
    \begin{subfigure}{0.5\textwidth}
        \includegraphics[width=\linewidth]{nr_SFO_0.445.png}
        \caption{}
        \label{fig:sfo_0.445_nr}
    \end{subfigure}%
    \begin{subfigure}{0.5\textwidth}
        \includegraphics[width=\linewidth]{gamma_tau_SFO_0.445.png}
        \caption{}
        \label{fig:sfo_0.445_gamma_tau}
    \end{subfigure}
    \caption{Here are represented the important features about the phase transition in San Francisco for UCI = 0.445.
    In \ref{fig:sfo_0.445_nt} is represented the number of people trapped in the road network at time t. In \ref{fig:sfo_0.445_scaled_nt}, the vertical line representing $t_{\infty}$. we see the values $n(t)t^{\alpha}$ over $\epsilon (t-t_0)$, with $\epsilon = \frac{R^2 - R_c^2}{R_c^2}$ with t up to $t_{start} + t_{end_{fit}}$ In \ref{fig:sfo_0.445_nr} showing the phase transition requirement. We plot n(R,$t_{\infty}$) with $\beta = $ 2.152. In \ref{fig:sfo_0.445_gamma_tau} we see the linear increase of $\tau$ compared to $\Gamma$}
    \label{fig:phase_transition_sfo_0.445}
\end{figure}

\begin{figure}
    \centering
    \begin{subfigure}{0.5\textwidth}
        \includegraphics[width=\linewidth]{nt_SFO_0.463.png}
        \caption{}
        \label{fig:sfo_0.463_nt}
    \end{subfigure}%
    \begin{subfigure}{0.5\textwidth}
        \includegraphics[width=\linewidth]{scaled_nt_SFO_0.463.png}
        \caption{}
        \label{fig:sfo_0.463_scaled_nt}
    \end{subfigure}
    
    \begin{subfigure}{0.5\textwidth}
        \includegraphics[width=\linewidth]{nr_SFO_0.463.png}
        \caption{}
        \label{fig:sfo_0.463_nr}
    \end{subfigure}%
    \begin{subfigure}{0.5\textwidth}
        \includegraphics[width=\linewidth]{gamma_tau_SFO_0.463.png}
        \caption{}
        \label{fig:sfo_0.463_gamma_tau}
    \end{subfigure}
    \caption{Here are represented the important features about the phase transition in San Francisco for UCI = 0.463.
    In \ref{fig:sfo_0.463_nt} is represented the number of people trapped in the road network at time t. In \ref{fig:sfo_0.463_scaled_nt}, the vertical line representing $t_{\infty}$. we see the values $n(t)t^{\alpha}$ over $\epsilon (t-t_0)$, with $\epsilon = \frac{R^2 - R_c^2}{R_c^2}$ with t up to $t_{start} + t_{end_{fit}}$ In \ref{fig:sfo_0.463_nr} showing the phase transition requirement. We plot n(R,$t_{\infty}$) with $\beta = $ 1.598. In \ref{fig:sfo_0.463_gamma_tau} we see the linear increase of $\tau$ compared to $\Gamma$}
    \label{fig:phase_transition_sfo_0.463}
\end{figure}

\begin{figure}
    \centering
    \begin{subfigure}{0.5\textwidth}
        \includegraphics[width=\linewidth]{nt_SFO_0.468.png}
        \caption{}
        \label{fig:sfo_0.468_nt}
    \end{subfigure}%
    \begin{subfigure}{0.5\textwidth}
        \includegraphics[width=\linewidth]{scaled_nt_SFO_0.468.png}
        \caption{}
        \label{fig:sfo_0.468_scaled_nt}
    \end{subfigure}
    
    \begin{subfigure}{0.5\textwidth}
        \includegraphics[width=\linewidth]{nr_SFO_0.468.png}
        \caption{}
        \label{fig:sfo_0.468_nr}
    \end{subfigure}%
    \begin{subfigure}{0.5\textwidth}
        \includegraphics[width=\linewidth]{gamma_tau_SFO_0.468.png}
        \caption{}
        \label{fig:sfo_0.468_gamma_tau}
    \end{subfigure}
    \caption{Here are represented the important features about the phase transition in San Francisco for UCI = 0.468.
    In \ref{fig:sfo_0.468_nt} is represented the number of people trapped in the road network at time t. In \ref{fig:sfo_0.468_scaled_nt}, the vertical line representing $t_{\infty}$. we see the values $n(t)t^{\alpha}$ over $\epsilon (t-t_0)$, with $\epsilon = \frac{R^2 - R_c^2}{R_c^2}$ with t up to $t_{start} + t_{end_{fit}}$ In \ref{fig:sfo_0.468_nr} showing the phase transition requirement. We plot n(R,$t_{\infty}$) with $\beta = $ 1.787. In \ref{fig:sfo_0.468_gamma_tau} we see the linear increase of $\tau$ compared to $\Gamma$}
    \label{fig:phase_transition_sfo_0.468}
\end{figure}

\begin{figure}
    \centering
    \begin{subfigure}{0.5\textwidth}
        \includegraphics[width=\linewidth]{nt_SFO_0.469.png}
        \caption{}
        \label{fig:sfo_0.469_nt}
    \end{subfigure}%
    \begin{subfigure}{0.5\textwidth}
        \includegraphics[width=\linewidth]{scaled_nt_SFO_0.469.png}
        \caption{}
        \label{fig:sfo_0.469_scaled_nt}
    \end{subfigure}
    
    \begin{subfigure}{0.5\textwidth}
        \includegraphics[width=\linewidth]{nr_SFO_0.469.png}
        \caption{}
        \label{fig:sfo_0.469_nr}
    \end{subfigure}%
    \begin{subfigure}{0.5\textwidth}
        \includegraphics[width=\linewidth]{gamma_tau_SFO_0.469.png}
        \caption{}
        \label{fig:sfo_0.469_gamma_tau}
    \end{subfigure}
    \caption{Here are represented the important features about the phase transition in San Francisco for UCI = 0.469.
    In \ref{fig:sfo_0.469_nt} is represented the number of people trapped in the road network at time t. In \ref{fig:sfo_0.469_scaled_nt}, the vertical line representing $t_{\infty}$. we see the values $n(t)t^{\alpha}$ over $\epsilon (t-t_0)$, with $\epsilon = \frac{R^2 - R_c^2}{R_c^2}$ with t up to $t_{start} + t_{end_{fit}}$ In \ref{fig:sfo_0.469_nr} showing the phase transition requirement. We plot n(R,$t_{\infty}$) with $\beta = $ 1.708. In \ref{fig:sfo_0.469_gamma_tau} we see the linear increase of $\tau$ compared to $\Gamma$}
    \label{fig:phase_transition_sfo_0.469}
\end{figure}

\begin{figure}
    \centering
    \begin{subfigure}{0.5\textwidth}
        \includegraphics[width=\linewidth]{nt_SFO_0.515.png}
        \caption{}
        \label{fig:sfo_0.515_nt}
    \end{subfigure}%
    \begin{subfigure}{0.5\textwidth}
        \includegraphics[width=\linewidth]{scaled_nt_SFO_0.515.png}
        \caption{}
        \label{fig:sfo_0.515_scaled_nt}
    \end{subfigure}
    
    \begin{subfigure}{0.5\textwidth}
        \includegraphics[width=\linewidth]{nr_SFO_0.515.png}
        \caption{}
        \label{fig:sfo_0.515_nr}
    \end{subfigure}%
    \begin{subfigure}{0.5\textwidth}
        \includegraphics[width=\linewidth]{gamma_tau_SFO_0.515.png}
        \caption{}
        \label{fig:sfo_0.515_gamma_tau}
    \end{subfigure}
    \caption{Here are represented the important features about the phase transition in San Francisco for UCI = 0.515.
    In \ref{fig:sfo_0.515_nt} is represented the number of people trapped in the road network at time t. In \ref{fig:sfo_0.515_scaled_nt}, the vertical line representing $t_{\infty}$. we see the values $n(t)t^{\alpha}$ over $\epsilon (t-t_0)$, with $\epsilon = \frac{R^2 - R_c^2}{R_c^2}$ with t up to $t_{start} + t_{end_{fit}}$ In \ref{fig:sfo_0.515_nr} showing the phase transition requirement. We plot n(R,$t_{\infty}$) with $\beta = $ 0.693. In \ref{fig:sfo_0.515_gamma_tau} we see the linear increase of $\tau$ compared to $\Gamma$}
    \label{fig:phase_transition_sfo_0.515}
\end{figure}

\begin{figure}
    \centering
    \begin{subfigure}{0.5\textwidth}
        \includegraphics[width=\linewidth]{nt_SFO_0.537.png}
        \caption{}
        \label{fig:sfo_0.537_nt}
    \end{subfigure}%
    \begin{subfigure}{0.5\textwidth}
        \includegraphics[width=\linewidth]{scaled_nt_SFO_0.537.png}
        \caption{}
        \label{fig:sfo_0.537_scaled_nt}
    \end{subfigure}
    
    \begin{subfigure}{0.5\textwidth}
        \includegraphics[width=\linewidth]{nr_SFO_0.537.png}
        \caption{}
        \label{fig:sfo_0.537_nr}
    \end{subfigure}%
    \begin{subfigure}{0.5\textwidth}
        \includegraphics[width=\linewidth]{gamma_tau_SFO_0.537.png}
        \caption{}
        \label{fig:sfo_0.537_gamma_tau}
    \end{subfigure}
    \caption{Here are represented the important features about the phase transition in San Francisco for UCI = 0.537.
    In \ref{fig:sfo_0.537_nt} is represented the number of people trapped in the road network at time t. In \ref{fig:sfo_0.537_scaled_nt}, the vertical line representing $t_{\infty}$. we see the values $n(t)t^{\alpha}$ over $\epsilon (t-t_0)$, with $\epsilon = \frac{R^2 - R_c^2}{R_c^2}$ with t up to $t_{start} + t_{end_{fit}}$ In \ref{fig:sfo_0.537_nr} showing the phase transition requirement. We plot n(R,$t_{\infty}$) with $\beta = $ 0.344. In \ref{fig:sfo_0.537_gamma_tau} we see the linear increase of $\tau$ compared to $\Gamma$}
    \label{fig:phase_transition_sfo_0.537}
\end{figure}

\begin{figure}
    \centering
    \begin{subfigure}{0.5\textwidth}
        \includegraphics[width=\linewidth]{nt_SFO_0.564.png}
        \caption{}
        \label{fig:sfo_0.564_nt}
    \end{subfigure}%
    \begin{subfigure}{0.5\textwidth}
        \includegraphics[width=\linewidth]{scaled_nt_SFO_0.564.png}
        \caption{}
        \label{fig:sfo_0.564_scaled_nt}
    \end{subfigure}
    
    \begin{subfigure}{0.5\textwidth}
        \includegraphics[width=\linewidth]{nr_SFO_0.564.png}
        \caption{}
        \label{fig:sfo_0.564_nr}
    \end{subfigure}%
    \begin{subfigure}{0.5\textwidth}
        \includegraphics[width=\linewidth]{gamma_tau_SFO_0.564.png}
        \caption{}
        \label{fig:sfo_0.564_gamma_tau}
    \end{subfigure}
    \caption{Here are represented the important features about the phase transition in San Francisco for UCI = 0.564.
    In \ref{fig:sfo_0.564_nt} is represented the number of people trapped in the road network at time t. In \ref{fig:sfo_0.564_scaled_nt}, the vertical line representing $t_{\infty}$. we see the values $n(t)t^{\alpha}$ over $\epsilon (t-t_0)$, with $\epsilon = \frac{R^2 - R_c^2}{R_c^2}$ with t up to $t_{start} + t_{end_{fit}}$ In \ref{fig:sfo_0.564_nr} showing the phase transition requirement. We plot n(R,$t_{\infty}$) with $\beta = $ 0.299. In \ref{fig:sfo_0.564_gamma_tau} we see the linear increase of $\tau$ compared to $\Gamma$}
    \label{fig:phase_transition_sfo_0.564}
\end{figure}

\begin{figure}
    \centering
    \begin{subfigure}{0.5\textwidth}
        \includegraphics[width=\linewidth]{nt_SFO_0.566.png}
        \caption{}
        \label{fig:sfo_0.566_nt}
    \end{subfigure}%
    \begin{subfigure}{0.5\textwidth}
        \includegraphics[width=\linewidth]{scaled_nt_SFO_0.566.png}
        \caption{}
        \label{fig:sfo_0.566_scaled_nt}
    \end{subfigure}
    
    \begin{subfigure}{0.5\textwidth}
        \includegraphics[width=\linewidth]{nr_SFO_0.566.png}
        \caption{}
        \label{fig:sfo_0.566_nr}
    \end{subfigure}%
    \begin{subfigure}{0.5\textwidth}
        \includegraphics[width=\linewidth]{gamma_tau_SFO_0.566.png}
        \caption{}
        \label{fig:sfo_0.566_gamma_tau}
    \end{subfigure}
    \caption{Here are represented the important features about the phase transition in San Francisco for UCI = 0.566.
    In \ref{fig:sfo_0.566_nt} is represented the number of people trapped in the road network at time t. In \ref{fig:sfo_0.566_scaled_nt}, the vertical line representing $t_{\infty}$. we see the values $n(t)t^{\alpha}$ over $\epsilon (t-t_0)$, with $\epsilon = \frac{R^2 - R_c^2}{R_c^2}$ with t up to $t_{start} + t_{end_{fit}}$ In \ref{fig:sfo_0.566_nr} showing the phase transition requirement. We plot n(R,$t_{\infty}$) with $\beta = $ 0.946. In \ref{fig:sfo_0.566_gamma_tau} we see the linear increase of $\tau$ compared to $\Gamma$}
    \label{fig:phase_transition_sfo_0.566}
\end{figure}

\begin{figure}
    \centering
    \begin{subfigure}{0.5\textwidth}
        \includegraphics[width=\linewidth]{nt_SFO_0.595.png}
        \caption{}
        \label{fig:sfo_0.595_nt}
    \end{subfigure}%
    \begin{subfigure}{0.5\textwidth}
        \includegraphics[width=\linewidth]{scaled_nt_SFO_0.595.png}
        \caption{}
        \label{fig:sfo_0.595_scaled_nt}
    \end{subfigure}
    
    \begin{subfigure}{0.5\textwidth}
        \includegraphics[width=\linewidth]{nr_SFO_0.595.png}
        \caption{}
        \label{fig:sfo_0.595_nr}
    \end{subfigure}%
    \begin{subfigure}{0.5\textwidth}
        \includegraphics[width=\linewidth]{gamma_tau_SFO_0.595.png}
        \caption{}
        \label{fig:sfo_0.595_gamma_tau}
    \end{subfigure}
    \caption{Here are represented the important features about the phase transition in San Francisco for UCI = 0.595.
    In \ref{fig:sfo_0.595_nt} is represented the number of people trapped in the road network at time t. In \ref{fig:sfo_0.595_scaled_nt}, the vertical line representing $t_{\infty}$. we see the values $n(t)t^{\alpha}$ over $\epsilon (t-t_0)$, with $\epsilon = \frac{R^2 - R_c^2}{R_c^2}$ with t up to $t_{start} + t_{end_{fit}}$ In \ref{fig:sfo_0.595_nr} showing the phase transition requirement. We plot n(R,$t_{\infty}$) with $\beta = $ 0.545. In \ref{fig:sfo_0.595_gamma_tau} we see the linear increase of $\tau$ compared to $\Gamma$}
    \label{fig:phase_transition_sfo_0.595}
\end{figure}

\begin{figure}
    \centering
    \begin{subfigure}{0.5\textwidth}
        \includegraphics[width=\linewidth]{nt_SFO_0.603.png}
        \caption{}
        \label{fig:sfo_0.603_nt}
    \end{subfigure}%
    \begin{subfigure}{0.5\textwidth}
        \includegraphics[width=\linewidth]{scaled_nt_SFO_0.603.png}
        \caption{}
        \label{fig:sfo_0.603_scaled_nt}
    \end{subfigure}
    
    \begin{subfigure}{0.5\textwidth}
        \includegraphics[width=\linewidth]{nr_SFO_0.603.png}
        \caption{}
        \label{fig:sfo_0.603_nr}
    \end{subfigure}%
    \begin{subfigure}{0.5\textwidth}
        \includegraphics[width=\linewidth]{gamma_tau_SFO_0.603.png}
        \caption{}
        \label{fig:sfo_0.603_gamma_tau}
    \end{subfigure}
    \caption{Here are represented the important features about the phase transition in San Francisco for UCI = 0.603.
    In \ref{fig:sfo_0.603_nt} is represented the number of people trapped in the road network at time t. In \ref{fig:sfo_0.603_scaled_nt}, the vertical line representing $t_{\infty}$. we see the values $n(t)t^{\alpha}$ over $\epsilon (t-t_0)$, with $\epsilon = \frac{R^2 - R_c^2}{R_c^2}$ with t up to $t_{start} + t_{end_{fit}}$ In \ref{fig:sfo_0.603_nr} showing the phase transition requirement. We plot n(R,$t_{\infty}$) with $\beta = $ 0.367. In \ref{fig:sfo_0.603_gamma_tau} we see the linear increase of $\tau$ compared to $\Gamma$}
    \label{fig:phase_transition_sfo_0.603}
\end{figure}

\begin{figure}
    \centering
    \begin{subfigure}{0.5\textwidth}
        \includegraphics[width=\linewidth]{nt_SFO_0.607.png}
        \caption{}
        \label{fig:sfo_0.607_nt}
    \end{subfigure}%
    \begin{subfigure}{0.5\textwidth}
        \includegraphics[width=\linewidth]{scaled_nt_SFO_0.607.png}
        \caption{}
        \label{fig:sfo_0.607_scaled_nt}
    \end{subfigure}
    
    \begin{subfigure}{0.5\textwidth}
        \includegraphics[width=\linewidth]{nr_SFO_0.607.png}
        \caption{}
        \label{fig:sfo_0.607_nr}
    \end{subfigure}%
    \begin{subfigure}{0.5\textwidth}
        \includegraphics[width=\linewidth]{gamma_tau_SFO_0.607.png}
        \caption{}
        \label{fig:sfo_0.607_gamma_tau}
    \end{subfigure}
    \caption{Here are represented the important features about the phase transition in San Francisco for UCI = 0.607.
    In \ref{fig:sfo_0.607_nt} is represented the number of people trapped in the road network at time t. In \ref{fig:sfo_0.607_scaled_nt}, the vertical line representing $t_{\infty}$. we see the values $n(t)t^{\alpha}$ over $\epsilon (t-t_0)$, with $\epsilon = \frac{R^2 - R_c^2}{R_c^2}$ with t up to $t_{start} + t_{end_{fit}}$ In \ref{fig:sfo_0.607_nr} showing the phase transition requirement. We plot n(R,$t_{\infty}$) with $\beta = $ 1.248. In \ref{fig:sfo_0.607_gamma_tau} we see the linear increase of $\tau$ compared to $\Gamma$}
    \label{fig:phase_transition_sfo_0.607}
\end{figure}

\begin{figure}
    \centering
    \begin{subfigure}{0.5\textwidth}
        \includegraphics[width=\linewidth]{nt_SFO_0.611.png}
        \caption{}
        \label{fig:sfo_0.611_nt}
    \end{subfigure}%
    \begin{subfigure}{0.5\textwidth}
        \includegraphics[width=\linewidth]{scaled_nt_SFO_0.611.png}
        \caption{}
        \label{fig:sfo_0.611_scaled_nt}
    \end{subfigure}
    
    \begin{subfigure}{0.5\textwidth}
        \includegraphics[width=\linewidth]{nr_SFO_0.611.png}
        \caption{}
        \label{fig:sfo_0.611_nr}
    \end{subfigure}%
    \begin{subfigure}{0.5\textwidth}
        \includegraphics[width=\linewidth]{gamma_tau_SFO_0.611.png}
        \caption{}
        \label{fig:sfo_0.611_gamma_tau}
    \end{subfigure}
    \caption{Here are represented the important features about the phase transition in San Francisco for UCI = 0.611.
    In \ref{fig:sfo_0.611_nt} is represented the number of people trapped in the road network at time t. In \ref{fig:sfo_0.611_scaled_nt}, the vertical line representing $t_{\infty}$. we see the values $n(t)t^{\alpha}$ over $\epsilon (t-t_0)$, with $\epsilon = \frac{R^2 - R_c^2}{R_c^2}$ with t up to $t_{start} + t_{end_{fit}}$ In \ref{fig:sfo_0.611_nr} showing the phase transition requirement. We plot n(R,$t_{\infty}$) with $\beta = $ 0.594. In \ref{fig:sfo_0.611_gamma_tau} we see the linear increase of $\tau$ compared to $\Gamma$}
    \label{fig:phase_transition_sfo_0.611}
\end{figure}

\begin{figure}
    \centering
    \begin{subfigure}{0.5\textwidth}
        \includegraphics[width=\linewidth]{nt_SFO_0.632.png}
        \caption{}
        \label{fig:sfo_0.632_nt}
    \end{subfigure}%
    \begin{subfigure}{0.5\textwidth}
        \includegraphics[width=\linewidth]{scaled_nt_SFO_0.632.png}
        \caption{}
        \label{fig:sfo_0.632_scaled_nt}
    \end{subfigure}
    
    \begin{subfigure}{0.5\textwidth}
        \includegraphics[width=\linewidth]{nr_SFO_0.632.png}
        \caption{}
        \label{fig:sfo_0.632_nr}
    \end{subfigure}%
    \begin{subfigure}{0.5\textwidth}
        \includegraphics[width=\linewidth]{gamma_tau_SFO_0.632.png}
        \caption{}
        \label{fig:sfo_0.632_gamma_tau}
    \end{subfigure}
    \caption{Here are represented the important features about the phase transition in San Francisco for UCI = 0.632.
    In \ref{fig:sfo_0.632_nt} is represented the number of people trapped in the road network at time t. In \ref{fig:sfo_0.632_scaled_nt}, the vertical line representing $t_{\infty}$. we see the values $n(t)t^{\alpha}$ over $\epsilon (t-t_0)$, with $\epsilon = \frac{R^2 - R_c^2}{R_c^2}$ with t up to $t_{start} + t_{end_{fit}}$ In \ref{fig:sfo_0.632_nr} showing the phase transition requirement. We plot n(R,$t_{\infty}$) with $\beta = $ 0.700. In \ref{fig:sfo_0.632_gamma_tau} we see the linear increase of $\tau$ compared to $\Gamma$}
    \label{fig:phase_transition_sfo_0.632}
\end{figure}

\begin{figure}
    \centering
    \begin{subfigure}{0.5\textwidth}
        \includegraphics[width=\linewidth]{nt_SFO_0.7.png}
        \caption{}
        \label{fig:sfo_0.7_nt}
    \end{subfigure}%
    \begin{subfigure}{0.5\textwidth}
        \includegraphics[width=\linewidth]{scaled_nt_SFO_0.7.png}
        \caption{}
        \label{fig:sfo_0.7_scaled_nt}
    \end{subfigure}
    
    \begin{subfigure}{0.5\textwidth}
        \includegraphics[width=\linewidth]{nr_SFO_0.7.png}
        \caption{}
        \label{fig:sfo_0.7_nr}
    \end{subfigure}%
    \begin{subfigure}{0.5\textwidth}
        \includegraphics[width=\linewidth]{gamma_tau_SFO_0.7.png}
        \caption{}
        \label{fig:sfo_0.7_gamma_tau}
    \end{subfigure}
    \caption{Here are represented the important features about the phase transition in San Francisco for UCI = 0.7.
    In \ref{fig:sfo_0.7_nt} is represented the number of people trapped in the road network at time t. In \ref{fig:sfo_0.7_scaled_nt}, the vertical line representing $t_{\infty}$. we see the values $n(t)t^{\alpha}$ over $\epsilon (t-t_0)$, with $\epsilon = \frac{R^2 - R_c^2}{R_c^2}$ with t up to $t_{start} + t_{end_{fit}}$ In \ref{fig:sfo_0.7_nr} showing the phase transition requirement. We plot n(R,$t_{\infty}$) with $\beta = $ 0.198. In \ref{fig:sfo_0.7_gamma_tau} we see the linear increase of $\tau$ compared to $\Gamma$}
    \label{fig:phase_transition_sfo_0.7}
\end{figure}

\begin{figure}
    \centering
    \begin{subfigure}{0.5\textwidth}
        \includegraphics[width=\linewidth]{nt_SFO_0.703.png}
        \caption{}
        \label{fig:sfo_0.703_nt}
    \end{subfigure}%
    \begin{subfigure}{0.5\textwidth}
        \includegraphics[width=\linewidth]{scaled_nt_SFO_0.703.png}
        \caption{}
        \label{fig:sfo_0.703_scaled_nt}
    \end{subfigure}
    
    \begin{subfigure}{0.5\textwidth}
        \includegraphics[width=\linewidth]{nr_SFO_0.703.png}
        \caption{}
        \label{fig:sfo_0.703_nr}
    \end{subfigure}%
    \begin{subfigure}{0.5\textwidth}
        \includegraphics[width=\linewidth]{gamma_tau_SFO_0.703.png}
        \caption{}
        \label{fig:sfo_0.703_gamma_tau}
    \end{subfigure}
    \caption{Here are represented the important features about the phase transition in San Francisco for UCI = 0.703.
    In \ref{fig:sfo_0.703_nt} is represented the number of people trapped in the road network at time t. In \ref{fig:sfo_0.703_scaled_nt}, the vertical line representing $t_{\infty}$. we see the values $n(t)t^{\alpha}$ over $\epsilon (t-t_0)$, with $\epsilon = \frac{R^2 - R_c^2}{R_c^2}$ with t up to $t_{start} + t_{end_{fit}}$ In \ref{fig:sfo_0.703_nr} showing the phase transition requirement. We plot n(R,$t_{\infty}$) with $\beta = $ 0.077. In \ref{fig:sfo_0.703_gamma_tau} we see the linear increase of $\tau$ compared to $\Gamma$}
    \label{fig:phase_transition_sfo_0.703}
\end{figure}

\begin{figure}
    \centering
    \begin{subfigure}{0.5\textwidth}
        \includegraphics[width=\linewidth]{nt_SFO_0.744.png}
        \caption{}
        \label{fig:sfo_0.744_nt}
    \end{subfigure}%
    \begin{subfigure}{0.5\textwidth}
        \includegraphics[width=\linewidth]{scaled_nt_SFO_0.744.png}
        \caption{}
        \label{fig:sfo_0.744_scaled_nt}
    \end{subfigure}
    
    \begin{subfigure}{0.5\textwidth}
        \includegraphics[width=\linewidth]{nr_SFO_0.744.png}
        \caption{}
        \label{fig:sfo_0.744_nr}
    \end{subfigure}%
    \begin{subfigure}{0.5\textwidth}
        \includegraphics[width=\linewidth]{gamma_tau_SFO_0.744.png}
        \caption{}
        \label{fig:sfo_0.744_gamma_tau}
    \end{subfigure}
    \caption{Here are represented the important features about the phase transition in San Francisco for UCI = 0.744.
    In \ref{fig:sfo_0.744_nt} is represented the number of people trapped in the road network at time t. In \ref{fig:sfo_0.744_scaled_nt}, the vertical line representing $t_{\infty}$. we see the values $n(t)t^{\alpha}$ over $\epsilon (t-t_0)$, with $\epsilon = \frac{R^2 - R_c^2}{R_c^2}$ with t up to $t_{start} + t_{end_{fit}}$ In \ref{fig:sfo_0.744_nr} showing the phase transition requirement. We plot n(R,$t_{\infty}$) with $\beta = $ 0.834. In \ref{fig:sfo_0.744_gamma_tau} we see the linear increase of $\tau$ compared to $\Gamma$}
    \label{fig:phase_transition_sfo_0.744}
\end{figure}

\begin{figure}
    \centering
    \begin{subfigure}{0.5\textwidth}
        \includegraphics[width=\linewidth]{nt_SFO_0.75.png}
        \caption{}
        \label{fig:sfo_0.75_nt}
    \end{subfigure}%
    \begin{subfigure}{0.5\textwidth}
        \includegraphics[width=\linewidth]{scaled_nt_SFO_0.75.png}
        \caption{}
        \label{fig:sfo_0.75_scaled_nt}
    \end{subfigure}
    
    \begin{subfigure}{0.5\textwidth}
        \includegraphics[width=\linewidth]{nr_SFO_0.75.png}
        \caption{}
        \label{fig:sfo_0.75_nr}
    \end{subfigure}%
    \begin{subfigure}{0.5\textwidth}
        \includegraphics[width=\linewidth]{gamma_tau_SFO_0.75.png}
        \caption{}
        \label{fig:sfo_0.75_gamma_tau}
    \end{subfigure}
    \caption{Here are represented the important features about the phase transition in San Francisco for UCI = 0.75.
    In \ref{fig:sfo_0.75_nt} is represented the number of people trapped in the road network at time t. In \ref{fig:sfo_0.75_scaled_nt}, the vertical line representing $t_{\infty}$. we see the values $n(t)t^{\alpha}$ over $\epsilon (t-t_0)$, with $\epsilon = \frac{R^2 - R_c^2}{R_c^2}$ with t up to $t_{start} + t_{end_{fit}}$ In \ref{fig:sfo_0.75_nr} showing the phase transition requirement. We plot n(R,$t_{\infty}$) with $\beta = $ 1.013. In \ref{fig:sfo_0.75_gamma_tau} we see the linear increase of $\tau$ compared to $\Gamma$}
    \label{fig:phase_transition_sfo_0.75}
\end{figure}

\begin{figure}
    \centering
    \begin{subfigure}{0.5\textwidth}
        \includegraphics[width=\linewidth]{nt_SFO_0.832.png}
        \caption{}
        \label{fig:sfo_0.832_nt}
    \end{subfigure}%
    \begin{subfigure}{0.5\textwidth}
        \includegraphics[width=\linewidth]{scaled_nt_SFO_0.832.png}
        \caption{}
        \label{fig:sfo_0.832_scaled_nt}
    \end{subfigure}
    
    \begin{subfigure}{0.5\textwidth}
        \includegraphics[width=\linewidth]{nr_SFO_0.832.png}
        \caption{}
        \label{fig:sfo_0.832_nr}
    \end{subfigure}%
    \begin{subfigure}{0.5\textwidth}
        \includegraphics[width=\linewidth]{gamma_tau_SFO_0.832.png}
        \caption{}
        \label{fig:sfo_0.832_gamma_tau}
    \end{subfigure}
    \caption{Here are represented the important features about the phase transition in San Francisco for UCI = 0.832.
    In \ref{fig:sfo_0.832_nt} is represented the number of people trapped in the road network at time t. In \ref{fig:sfo_0.832_scaled_nt}, the vertical line representing $t_{\infty}$. we see the values $n(t)t^{\alpha}$ over $\epsilon (t-t_0)$, with $\epsilon = \frac{R^2 - R_c^2}{R_c^2}$ with t up to $t_{start} + t_{end_{fit}}$ In \ref{fig:sfo_0.832_nr} showing the phase transition requirement. We plot n(R,$t_{\infty}$) with $\beta = $ 0.597. In \ref{fig:sfo_0.832_gamma_tau} we see the linear increase of $\tau$ compared to $\Gamma$}
    \label{fig:phase_transition_sfo_0.832}
\end{figure}

\begin{figure}
    \centering
    \begin{subfigure}{0.5\textwidth}
        \includegraphics[width=\linewidth]{nt_SFO_0.841.png}
        \caption{}
        \label{fig:sfo_0.841_nt}
    \end{subfigure}%
    \begin{subfigure}{0.5\textwidth}
        \includegraphics[width=\linewidth]{scaled_nt_SFO_0.841.png}
        \caption{}
        \label{fig:sfo_0.841_scaled_nt}
    \end{subfigure}
    
    \begin{subfigure}{0.5\textwidth}
        \includegraphics[width=\linewidth]{nr_SFO_0.841.png}
        \caption{}
        \label{fig:sfo_0.841_nr}
    \end{subfigure}%
    \begin{subfigure}{0.5\textwidth}
        \includegraphics[width=\linewidth]{gamma_tau_SFO_0.841.png}
        \caption{}
        \label{fig:sfo_0.841_gamma_tau}
    \end{subfigure}
    \caption{Here are represented the important features about the phase transition in San Francisco for UCI = 0.841.
    In \ref{fig:sfo_0.841_nt} is represented the number of people trapped in the road network at time t. In \ref{fig:sfo_0.841_scaled_nt}, the vertical line representing $t_{\infty}$. we see the values $n(t)t^{\alpha}$ over $\epsilon (t-t_0)$, with $\epsilon = \frac{R^2 - R_c^2}{R_c^2}$ with t up to $t_{start} + t_{end_{fit}}$ In \ref{fig:sfo_0.841_nr} showing the phase transition requirement. We plot n(R,$t_{\infty}$) with $\beta = $ 0.631. In \ref{fig:sfo_0.841_gamma_tau} we see the linear increase of $\tau$ compared to $\Gamma$}
    \label{fig:phase_transition_sfo_0.841}
\end{figure}

